\section*{Bab \ref{ch:g5:food-webs} --- Jaring Makanan}

\begin{solution}{exo:g5:food-webs:1}
Jaringan semua \omterm{def:g3:food-chains:chain}{rantai makanan} sebuah \omterm{def:g2:where-animals-live:habitat}{habitat} yang digambar sekaligus
--- setiap \omterm{def:g5:biodiversity:species}{spesies} satu kali, satu anak panah untuk setiap hubungan
santapan. Rantainya bersilangan karena kebanyakan \omterm{def:g1:animals-around-us:animal}{hewan} memakan
beberapa \omterm{def:g5:biodiversity:species}{spesies} dan \omterm{def:g3:food-chains:chain}{dimakan oleh} beberapa \omterm{def:g5:biodiversity:species}{spesies}.
\end{solution}

\begin{solution}{exo:g5:food-webs:2}
\omterm{def:g5:food-webs:roles}{Produsen}: \omterm{def:g1:plants-around-us:plant}{tumbuhan} (rumput, semanggi, \omterm{def:g1:plants-around-us:parts}{daun} pohon ek). \omterm{def:g5:food-webs:roles}{Konsumen}: \omterm{def:g1:animals-around-us:animal}{hewan}
--- \omterm{def:g2:what-animals-eat:kinds}{herbivora} seperti kelinci, \omterm{def:g2:what-animals-eat:kinds}{karnivora} seperti elang. \omterm{def:g5:food-webs:roles}{Pengurai}: para
pelapuk yang makan dari sisa yang mati (tidak ada yang digambar pada
gambarnya --- mereka bekerja di balik bayang-bayang di bawah semuanya).
\end{solution}

\begin{solution}{exo:g5:food-webs:3}
Gelatik birunya memakan \omterm{ex:g2:animals-grow-up:butterfly}{ulat} dan belalang; ia \omterm{def:g3:food-chains:chain}{dimakan oleh} alap-alap.
\end{solution}

\begin{solution}{exo:g5:food-webs:4}
Rumput $\rightarrow$ belalang $\rightarrow$ kadal
$\rightarrow$ elang, dan rumput $\rightarrow$ belalang
$\rightarrow$ gelatik biru $\rightarrow$ alap-alap.
\end{solution}

\begin{solution}{exo:g5:food-webs:5}
\omterm{def:g5:food-webs:roles}{Produsen} banyak, \omterm{def:g2:what-animals-eat:kinds}{herbivora} lebih sedikit, \omterm{def:g2:what-animals-eat:kinds}{karnivora} sedikit. Karena
tiap lantai makan dari lantai di bawahnya: pemakan tidak akan pernah
dapat melebihi jumlah yang tersedia untuk dimakan.
\end{solution}

\begin{solution}{exo:g5:food-webs:6}
Karena sepasang elang membutuhkan belalang senilai satu padang rumput,
yang berlapis beberapa kadal, untuk hidup --- puncak piramidanya
berdiri di atas lantai-lantai lebar di bawahnya.
\end{solution}

\begin{solution}{exo:g5:food-webs:7}
Pemburunya --- rubah, elang --- makan lebih baik lalu berlipat ganda
sesudahnya; perburuan tambahan itu melelehkan kelebihan kelincinya,
dan bilangannya mengendap kembali menuju keseimbangan.
\end{solution}

\begin{solution}{exo:g5:food-webs:8}
Jawaban terbuka. Sebuah jaring kebun yang dapat dipertahankan alasannya:
selada dan rumput (\omterm{def:g5:food-webs:roles}{produsen}) $\rightarrow$ siput, kutu \omterm{def:g1:plants-around-us:parts}{daun}, \omterm{prop:g3:sorting-living-things:groups}{burung}
pipit; siput $\rightarrow$ landak; kutu \omterm{def:g1:plants-around-us:parts}{daun} $\rightarrow$ kepik;
\omterm{prop:g3:sorting-living-things:groups}{burung} pipit dan kepik $\rightarrow$ kucing tetangga atau alap-alap.
\end{solution}

\begin{solution}{exo:g5:food-webs:9}
Langkah pertama: belalang, yang lebih sedikit dimakan, bertambah;
elangnya, yang menjadi lebih miskin sebanyak satu mangsa, bersandar
lebih keras pada kelinci (anak panah cadangannya). Langkah kedua:
belalang yang lebih banyak menekan rumputnya lebih keras; perburuan
kelinci yang lebih banyak menyorong kelincinya turun --- yang pada
gilirannya meringankan rumputnya. Jaringnya melengkung lalu sebagian
menyerap kehilangannya; \omterm{def:g5:biodiversity:species}{spesies} yang tidak punya anak panah cadangan
akan paling menderita, dan di sini semuanya punya cadangan.
\end{solution}

\begin{solution}{exo:g5:food-webs:10}
Elangnya disingkirkan; lalu tikus dan tikus sawah, yang terbebas dari
pemburu utamanya, berlipat ganda; lalu \omterm{def:g1:animals-around-us:animal}{hewan} pengerat yang berlipat
ganda itu memakan panennya --- gandum yang lebih banyak daripada yang
pernah dibebankan elangnya. Para pengampanyenya meramalkan langkah
pertama lalu berhenti; jaringnya berjalan dua langkah.
\end{solution}

\begin{solution}{exo:g5:food-webs:11}
Kehilangan pohon eknya memukul yang bersandar padanya ke arah atas:
\omterm{ex:g2:animals-grow-up:butterfly}{ulatnya} (\omterm{def:g1:plants-around-us:parts}{daun} pohon ek adalah satu-satunya anak panahnya) kelaparan,
lalu gelatik birunya bersandar sepenuhnya pada belalang, lalu
alap-alapnya merasakan penipisannya --- sebuah kemarau dari bawah ke
atas seperti musim panas yang kering. Kehilangan alap-alapnya merambat
ke bawah: gelatik biru, yang lebih sedikit diburu, bertambah lalu
menekan \omterm{ex:g2:animals-grow-up:butterfly}{ulat} dan belalangnya --- kisah penyingkiran pemburu. Di sini
kehilangan alap-alapnya lebih mudah diserap: mangsanya punya pemburu
lain dan gangguannya memudar lantai demi lantai, sedangkan yang
bersandar pada pohon eknya sama sekali tidak punya anak panah cadangan
--- seluruh lantai \omterm{ex:g2:animals-grow-up:butterfly}{ulatnya} berdiri di atas satu \omterm{def:g5:food-webs:roles}{produsen}.
\end{solution}
